% GENERATED FILE. Edit canonical lessons, then run npm run generate:books.
% canonical-source-hash: fnv1a64:dc806a11382e6820
% canonical-lessons: GE-C257-hin-und-zurueck, GE-C257-zug-verpasst, GE-C257-review

\chapter{At the Station — hin und zurück, den Zug verpasst}
\label{ch:ge-at-the-station}
\hlchaptermodality{257}

\begin{chapteropening}
\textbf{By the end of this chapter:} \emph{I can buy a ticket one way or there and back, and report a problem such as a missed train.}

Complete GE-C257-review: buy a return ticket to a city, then report a missed train and ask what to do next.
\end{chapteropening}

\section[hin und zurück]{hin und zurück — there and back}

\label{lesson:GE-C257-hin-und-zurueck}



\begin{quote}
Say \textquotedblleft{}the ticket\textquotedblright{}. (\emph{Die Fahrkarte.}) And \textquotedblleft{}simple, easy\textquotedblright{}. (\emph{Einfach.}) At a ticket counter, \emph{einfach} has a second job.
\end{quote}

\subsection*{You'll want to know: hin und zurück}
\textbf{hin und zurück} — \textquotedblleft{}there and back\textquotedblright{}, a return: \emph{hin oont tsoo-RÜCK}.

The counter question is short, and so is the answer:

\begin{quote}
\emph{Einfach oder hin und zurück?} — One way or return? \emph{Hin und zurück, bitte.} — A return, please. \emph{Was kostet das?} — How much is that?
\end{quote}

\textbf{Einfach} at the counter means a single, one way. \textbf{Hin} is \textquotedblleft{}there, away from here\textquotedblright{}, and \textbf{zurück} is \textquotedblleft{}back\textquotedblright{}. The return half of the journey is the word you already have: \emph{die Rückfahrt}.

\begin{grammarlens}[title={Grammar Lens: nach for a city, zum for a place in town}]
A ticket goes \textbf{nach} a city, or a country that takes no article:

\begin{quote}
\emph{Eine Fahrkarte nach Berlin, bitte.} — A ticket to Berlin, please. \emph{Einmal nach Hamburg, hin und zurück.} — One return to Hamburg.
\end{quote}

Inside the town, you went \emph{zum Bahnhof} and \emph{zur Post}. Between towns, you travel \emph{nach Berlin}.

Then the two questions for the platform:

\begin{quote}
\emph{Wann fährt der Zug ab?} — When does the train leave? \emph{Von welchem Gleis?} — From which platform?
\end{quote}

In Switzerland the ticket is often \emph{das Billett}, and a return is \emph{retour}: \emph{einmal Bern retour}. \emph{Fahrkarte} is understood everywhere.
\end{grammarlens}

\begin{cousinweb}
\textbf{zurück} is usually explained as \emph{zu}, \textquotedblleft{}to\textquotedblright{}, + \emph{Rück-}, \textquotedblleft{}back\textquotedblright{}: in Middle High German, \emph{ze rucke}, \textquotedblleft{}to the back\textquotedblright{}. The \emph{Rück-} is the same as in \emph{der Rücken} and \emph{die Rückfahrt}; EtymDB 2.1 lists \emph{zurück} beside Dutch \emph{terug}, built the same way, but does not take it apart.

\textbf{hin} is listed by EtymDB beside English \emph{hence}, \textquotedblleft{}from here\textquotedblright{}; both point away from where the speaker stands.
\end{cousinweb}

\subsection*{Your turn}
Say these aloud:

\begin{itemize}
\raggedright
  \item \emph{hin und zurück}
  \item \emph{Eine Fahrkarte nach Berlin, bitte.}
  \item the clerk: \emph{Einfach oder hin und zurück?}
  \item you: \emph{Hin und zurück, bitte. Was kostet das?}
  \item \emph{Wann fährt der Zug ab? Von welchem Gleis?}
\end{itemize}

\begin{tcolorbox}[breakable,colback=teal!4,colframe=teal!35!black,title={Before you move on}]
\textquotedblleft{}A return, please\textquotedblright{}? (\emph{\textbf{Hin und zurück, bitte.}}) \textquotedblleft{}One way\textquotedblright{}? (\emph{\textbf{Einfach.}}) \textquotedblleft{}To Berlin\textquotedblright{}: \emph{zum}, \emph{zur} or \emph{nach}? (\emph{\textbf{Nach Berlin.}})
\end{tcolorbox}
\section[ich habe den Zug verpasst]{ich habe den Zug verpasst — when something goes wrong}

\label{lesson:GE-C257-zug-verpasst}



\begin{quote}
Say \textquotedblleft{}to miss\textquotedblright{}, a train or a chance. (\emph{Verpassen.}) Then \textquotedblleft{}I said\textquotedblright{}, in the Perfekt. (\emph{Ich habe gesagt.})
\end{quote}

\subsection*{You'll want to know: ich habe den Zug verpasst}
\textbf{Ich habe den Zug verpasst.} — \textquotedblleft{}I missed the train\textquotedblright{}: \emph{ikh HAH-buh dehn tsook fer-PASST}.

Three things in it matter. It is the \textbf{Perfekt}, because a ticket counter is a conversation, and conversation mostly runs in the Perfekt. The participle is \textbf{verpasst}, with no \emph{ge-}: verbs that begin with \emph{ver-} do without it. And \emph{der Zug} becomes \textbf{den Zug}, because the train is the thing missed.

\begin{grammarlens}[title={Grammar Lens: say what happened, then ask}]
Reporting a problem goes in two moves. First the problem, in the Perfekt or the present:

\begin{quote}
\emph{Ich habe den Zug verpasst.} \emph{Ich habe meine Fahrkarte verloren.} — I've lost my ticket. \emph{Mein Koffer ist kaputt.} — My suitcase is broken.
\end{quote}

Then the question that asks for help:

\begin{quote}
\emph{Wann fährt der nächste Zug?} — When is the next train? \emph{Was kann ich machen?} — What can I do?
\end{quote}

\emph{Der nächste Zug} is the \emph{nächste} of \emph{nächste Woche}, next.
\end{grammarlens}

\begin{cousinweb}
\textbf{der Zug} $\leftarrow$ Middle High German \emph{zuc} (EtymDB 2.1), which EtymDB traces, beside Old English \emph{tyge}, to a Proto-Germanic word for \textquotedblleft{}a tug, a pull\textquotedblright{}. It belongs to the family of \textbf{ziehen}, \textquotedblleft{}to pull\textquotedblright{}: a train is the thing the engine pulls.

\textbf{verpassen}, from chapter one hundred and ninety, is \emph{ver-} + \emph{passen}, from Dutch, and further back from French \emph{passer}, \textquotedblleft{}to pass\textquotedblright{}: to let a thing pass you by.
\end{cousinweb}

\subsection*{Your turn}
Say these aloud:

\begin{itemize}
\raggedright
  \item \emph{Ich habe den Zug verpasst.}
  \item \emph{Wann fährt der nächste Zug?}
  \item \emph{Ich habe meine Fahrkarte verloren. Was kann ich machen?}
  \item \emph{Mein Koffer ist kaputt.}
  \item the whole problem: \emph{Ich habe den Zug nach Berlin verpasst. Wann fährt der nächste Zug?}
\end{itemize}

\begin{tcolorbox}[breakable,colback=teal!4,colframe=teal!35!black,title={Before you move on}]
\textquotedblleft{}I missed the train\textquotedblright{}? (\emph{\textbf{Ich habe den Zug verpasst.}}) Why no \emph{ge-} in \emph{verpasst}? (\textbf{It starts with} \emph{\textbf{ver-}}.) Why \emph{den Zug}? (\textbf{The train is the thing missed.}) \textquotedblleft{}When is the next train?\textquotedblright{} (\emph{\textbf{Wann fährt der nächste Zug?}})
\end{tcolorbox}
\section[(review)]{Review — a day at the station}

\label{lesson:GE-C257-review}



\begin{quote}
Ask a stranger the way to the station. (\emph{Entschuldigung, wie komme ich zum Bahnhof?}) And give the answer back. (\emph{Gehen Sie geradeaus.})
\end{quote}

\begin{grammarlens}[title={What you've built}]
Two chapters, and a trip by train is something you can manage in German. You find the station, you buy \emph{eine Fahrkarte nach Berlin}, \emph{einfach} or \emph{hin und zurück}, and when the day goes wrong you say what happened in the Perfekt and ask what to do.

Here it is in one go:

\begin{quote}
\emph{Einmal nach Berlin, hin und zurück, bitte. Was kostet das?} \emph{Ich habe den Zug verpasst. Wann fährt der nächste Zug?}
\end{quote}
\end{grammarlens}

\subsection*{Your turn}
Say these aloud:

\begin{itemize}
\raggedright
  \item \emph{Eine Fahrkarte nach Hamburg, bitte.}
  \item the clerk: \emph{Einfach oder hin und zurück?}
  \item you: \emph{Einfach, bitte. Was kostet das?}
  \item later: \emph{Ich habe den Zug verpasst. Was kann ich machen?}
  \item one more problem of your own, and one question about it
  \item \emph{Können Sie mir erklären, warum der Zug Verspätung hat?} Then the three steps: \emph{Ich komme heute spät. Ich komme spät, weil ich den Zug verpasst habe. Darum rufe ich dich an.} And advice: \emph{Man sollte den Fahrplan lesen.}
\end{itemize}

\begin{tcolorbox}[breakable,colback=teal!4,colframe=teal!35!black,title={Before you move on}]
\textquotedblleft{}A return\textquotedblright{}? (\emph{\textbf{Hin und zurück.}}) \textquotedblleft{}To Hamburg\textquotedblright{}? (\emph{\textbf{Nach Hamburg.}}) \textquotedblleft{}I missed the train\textquotedblright{}? (\emph{\textbf{Ich habe den Zug verpasst.}})
\end{tcolorbox}
